\section{Proofs}
\label{app:proof}
\subsection{Proof of \cref{lem:sim_impossible}}
\label{app:sim_impossible_proof}
\begin{proof}
    Algebraically, it is straightforward.
    Any image of abelian Lie algebra under Lie algebra homomorphism is abelian.
    There is no representation of abelian Lie algebra that depends on path order.
    
    Constructively, general SSM $\mbf{S}$ and abelian SSM $\hbf{S}$ are linear ODEs.
    It suffices to find the existence of nonsingular linear state transformation map $\mbf{T}$
    that satisfies $\mbf{T}\hat{h}(t) = h(t)$ for any time $t$ and admissible inputs
    (\cref{par:sim_min_main} and \cref{app:sim_detail}).

    Let us denote $(\Phi, \hat{\Phi})$ be the state-transition matrix of $(\mbf{S}, \hbf{S})$.
    It is clear that $\mbf{T}\hat{h}(t) = h(t)$ implies $\mbf{T}\hat{\Phi}(t,0) = \Phi(t,0)\mbf{T}$.
    Consider an admissible input path $x: [0,T] \to \mcal{X}$,
    and its measure-preserving permutation path $x'(t)= x(\sigma(t))$
    where $\sigma:[0,T] \to [0,T]$.
    Assume $x'$ is also admissible.
    Then, being abelian:
    \begin{equation*}
        \hat{\Phi}_{x}(T,0) = \hat{\Phi}_{x'}(T,0),
        \quad
        \mbf{T}(\hat{\Phi}_{x}-\hat{\Phi}_{x'})\mbf{T}^{-1} = \mbf{0}.
    \end{equation*}
    However, for non-abelian $\mbf{S}$, chronological exponential map depends on the path order and $\Phi_{x}(T,0) \neq \Phi_{x'}(T,0)$,
    contradicting the existence of $\mbf{T}$.
    
    As noted in \cref{app:sim_detail}, the proof generalizes to the case $m = \dim \hbf{S} \ge \dim \mbf{S} = n$.
    $\mbf{T}$ simply needs to satisfy:
    \begin{equation*}
        \mbf{T}\hat{h} = \begin{bmatrix}
            h \\
            *
        \end{bmatrix}
        \begin{matrix*}[l]
            \rbrace \; n \\
            \rbrace \; m-n,
        \end{matrix*}
    \end{equation*}
    for some variables $*$ to match dimension \cite{toth2010modeling}.
    The algebraic nature is intact, and the proof immediately follows.
    With respect to \cref{eq:SSM_sim_error},
    truncating $*$-corresponding rows of $\mbf{T}$ recover full-rank matrix $\mbf{P}$.
\end{proof}
\subsection{Proof of \cref{thm:sim_error}}
\label{app:sim_error_proof}
Here we divide the problem into abelian and restricted cases.
Notation follows \cref{app:sim_impossible_proof}.
\begin{proof}[Abelian SSM]
    Consider general SSM $\mbf{S}$.
    Assume an input path $x$ and permuted $x'$ on window $[0,T]$ are both admissible.
    With small enough $T$, we can derive a loose lower bound of $\|\Phi_{x}(T) - \Phi_{x'}(T)\|$ with Magnus expansion in \cref{app:CDE}:
    \begin{align*}
        \|\Phi_x(T) - \Phi_{x'}(T)\|
        &= \|\exp \Omega(T) - \exp \Omega'(T) \| \\
        & \ge \exp \Bigl(
        - \max
        \{ \|\Omega(T)\|, \| \Omega'(T) \| \}
        \Bigr)
        \| \Omega(T) - \Omega'(T) \| \\
        &= \lambda\| \Omega(T) - \Omega'(T) \|, \qquad \lambda \in (0,1].
    \end{align*}
    Recall \cref{eq:omega_2},
    \begin{align*}
        \|\Omega_2(T) - \Omega_2'(T)\|
        &= 
        \biggl\|
        \frac{1}{4}\iint_{[0,T]^2}
        \Bigl(
        \text{sgn}(t_1 - t_2)
        -
        \text{sgn}(\sigma(t_1) - \sigma(t_2))
        \Bigr)
        [\mbf{A}(x(t_1)), \mbf{A}(x(t_2))]
        \, dt_1 dt_2
        \biggr\| \\
        &\le 
        \biggl \|
        \frac{1}{2} \iint_{[0,T]^2}
        \text{sgn}(t_1 - t_2)
        [\mbf{A}(x(t_1)), \mbf{A}(x(t_2))]
        \, dt_1 dt_2
        \biggr \| \\
        &= 2\|\Omega_2(T)\|,
    \end{align*}
    where equality holds when $\sigma(t)$ is a \emph{time-reversal} function.

    Time-reversal function simplifies Magnus series further.
    It is straightforward that $\Omega(T) - \Omega'(T) = \sum_{i=1} \Omega_{2n}(T)$ with $\Omega_{2n}(T) = -\Omega_{2n}'(T)$ and $\Omega_{2n+1}(T) = \Omega_{2n+1}'(T)$.
    More generally,
    $\Omega_2$ is the leading term in $\Omega_{\ge 2}(T)$
    without assuming that the time-reversal function is a permutation.
    Its tail $\Omega_{\ge 3}(T)$ shrinks at $\mathcal{O}(T^3)$.
    Then, we can derive the rest of the lower bound:
    \begin{equation}
    \label{app:eq_OmegaError_2nd}
        \| \Omega(T) - \Omega'(T) \|
        =
        2\|\Omega_{2n}(T)\|
        \ge 2c\|\Omega_2(T)\|, \qquad c \in (0,1].
    \end{equation}
    On the other hand, \cref{app:sim_impossible_proof} implies abelian $\hbf{S}$ generates the same final state under path permutation.
    Thus, either $x$ or $x'$ should incur state mismatch:
    \begin{equation*}
        \max_{x,x'} \Delta_h
        \ge
        \lambda c \|\Omega_2(T)\|\|h_0\|,
    \end{equation*}
    provides lower bound of homogeneous simulation error incurred on small window.
    For anyone interested in output error,
    the minimality condition guarantees existence of some common suffix sequence
    that generates output mismatch from state mismatch.
    
    When $T$ is large, we decompose the path by small enough time partitions, as noted in \cref{eq:flow_composition}:
    \begin{equation}
        \Phi(T,0) = \Phi(T,t_k)\Phi(t_k,t_{k-1})\cdots\Phi(t_1,0).
    \end{equation}
    Note that \cref{eq:flow_composition} admits simple rearrange of terms:
    \begin{align*}
        \Phi_x(t_2,t_1)\Phi_x(t_1,0) - \Phi_{x'}(t_2,t_1)\Phi_{x'}(t_1,0)
        =
        \Phi_x(t_2,t_1)\bigl(\Phi_x(t_1,0) - \Phi_{x'}(t_1,0)\bigr)
        +
        \bigl(\Phi_x(t_2,t_1) - \Phi_{x'}(t_2,t_1)\bigr)\Phi_{x'}(t_1,0).
    \end{align*}
    Therefore,
    \begin{equation}
        \Phi_x(T) - \Phi_{x'}(T) = \sum_{i=1}^k \Phi_x(T, t_{i+1})
        \Bigl(
        \Phi_x(t_i, t_{i-1}) - \Phi_{x'}(t_i, t_{i-1})
        \Bigr)
        \Phi_{x'}(t_{i-1},0),
    \end{equation}
    and with minimality condition, $x$ can be designed to accumulate local error over each small pieces.
\end{proof}
\begin{proof}[Restricted SSM]
    Let us construct admissible inputs on interval $[0,2T]$ by composing two $T$-length sequence\footnote{
    Choice of $2T$ is arbitrary. Two sequences need not have the same length.
    }:
    \begin{equation*}
        x_1 = \text{prefix}_1 \oplus \mbf{xy},
        \quad
        x_1' = \text{prefix}_1 \oplus \mbf{yx},
        \quad
        x_2 = \text{prefix}_2 \oplus \mbf{xy},
        \quad
        x_2' = \text{prefix}_2 \oplus \mbf{yx}
    \end{equation*}
    with permutation map $\sigma$ on $(T, 2T]$.
    Here we abuse notation and denote $(\mbf{xy}, \mbf{yx})$ to emphasize permutation.
    
    Consider general SSM $\mbf{S}$.
    As $x_1$ and $x_1'$ shares prefix, $\Phi_{x_1}(T,0) = \Phi_{x_1'}(T,0)$.
    Shortly, we denote $\mbf{M} := \Phi_{x_1}(2T,T) - \Phi_{x_1'}(2T,T)$.
    By simple subtraction:
    \begin{align*}
        h_1(2T) - h_1'(2T)
        =
        \mbf{M}h_1(T)
        +
        \int_T^{2T}
        \Bigl(
        \Phi_{\mbf{xy}}(2T, \tau)\mbf{b}(\mbf{xy}(\tau))
        -
        \Phi_{\mbf{yx}}(2T, \tau)\mbf{b}(\mbf{yx}(\tau))
        \Bigr) d\tau.
    \end{align*}
    Similarly, we can obtain $h_2(2T) - h_2'(2T)$.
    Then, subtraction between terms cancels inhomogeneous part out:
    \begin{equation}
    \label{app:eq_subtraction_magic}
        \Bigl(
        h_1(2T) - h_1'(2T)
        \Bigr)
        -
        \Bigl(
        h_2(2T) - h_2'(2T)
        \Bigr)
        =
        \mbf{M}
        (h_1(T) - h_2(T)).
    \end{equation}
    On the other hand, for restricted $\hbf{S}$, $\mbf{M} = \mbf{0}$.
    Therefore, among these four inputs,
    \begin{align}
    \label{app:eq_inhomogeneous_error}
        \max_{x_1,x_1', x_2, x_2'} \Delta_h
        &\ge \frac{1}{4} \|\mbf{M}
        (h_1(T) - h_2(T))\|,
    \end{align}
    with magnitude of $\mbf{M}$ following \cref{app:eq_OmegaError_2nd}.
    
    With minimality condition,
    $\text{prefix}_1$ and $\text{prefix}_2$ can be chosen such that $h_1(T) \neq h_2(T)$ for $\mbf{S}$.
    Up to bounded magnitude, we can make $h_1(T)$ and $h_2(T)$ be as far as possible.
    Long sequence behavior is analogous to the abelian case.
\end{proof}
\subsection{Proof of \cref{lem:stackderivedlength}}
\label{app:stackderivedlength_proof}
We first show abelian $k$-layer SSM has up to $k$ derived length,
and naturally derives $2k$ for restricted SSMs.
To obtain notational simplicity, we choose vector field notation for this proof.
\begin{proof}[Abelian SSM]
    Consider abelian $k$-layer SSM with vector field
    $F(t,h,x) = (F_0, F_1, \ldots, F_{k-1})$ on a manifold $\mcal{M} \cong \bb{R}^n, \mcal{M}_i \cong \bb{R}^{n_i}$.
    Assume $(t,x)$ is fixed.
    For the layer index $i$ and state coordinate index $a$, we trivially extend the local vector field:
    \begin{equation}
    \label{app:eq_trivial_support_extension}
        h_i \in \mcal{M}_i,
        \quad
        \partial_{i,a} = \frac{\partial}{\partial h_i^a},
        \quad 
        \bar{F}_i(h) = \sum_{a=1}^{n_i} F_i^a (h_{\le i}) \partial_{i,a}
        = (\underbrace{0,\ldots,0}_{0:i-1}, F_i, \underbrace{0,\ldots,0}_{i+1:k-1}),
    \end{equation}
    such that the coefficient function $(\bar{F}_i)_j^c = 0$ for any $j \neq i$ and $\bar{F}_i^c = F_i^c$.
    Then the global vector field $F(h) = \sum_{i=0}^{k-1} \bar{F}_i(h)$.
    
    Now consider two fixed control inputs $x,y$ and corresponding vector fields $X,Y$
    such that $X(h) = F(h,x)$ and $Y(h) = F(h,y)$.
    Following \cref{app:eq_trivial_support_extension}, let us denote:
    \begin{align*}
        X = \sum_{i=0}^{k-1} X_i,
        \quad
        X_i &= \sum_{a=1}^{n_i}X_i^a \partial_{i,a},
    \end{align*}
    and similarly for $Y$.

    With \cref{app:eq_LieBracket_Directional} let us construct their Lie bracket $[X,Y] = \sum_{i,j}[X_i,Y_j]$.
    When $i=j$, trivially $[X_i,Y_i] = 0$ from abelian definition.
    For $i < j$, note that $X_i$ is independent of $h_j$ by construction, therefore:
    \begin{align}
        [X_i,Y_j]_i^c &= \sum_{a=1}^{n_i}X_i^a \partial_{i,a} (Y_j)_i^c
        - \sum_{b=1}^{n_j}Y_j^b \partial_{j,b} X_i^c = 0, \label{eq:zerocrossbracket} \\
        [X_i,Y_j]_j^c &= \sum_{a=1}^{n_i}X_i^a \partial_{i,a} Y_j^c
        - \sum_{b=1}^{n_j}Y_j^b \partial_{j,b} (X_i)_j^c
        = \sum_{a=1}^{n_i}X_i^a \partial_{i,a} Y_j^c, \label{eq:nonzerocrossbracket}
    \end{align}
    such that nonzero component of $[X_i,Y_j]$ only lives in the $j$-th layer.

    \cref{eq:nonzerocrossbracket} provides a chain of ideals.
    Let us denote $\fr{g}$ be the Lie algebra generated by the collection of global vector fields $F$.
    We define its ideal $I_{l} := \{X \in \fr{g} \mid X(h) = \sum_{i \ge l} X_i(h)\}$ with zero component below layer $l$.
    Then:
    \begin{equation*}
        [X_i,Y_j] \in I_{\max (i,j)},
        \quad
        [X,Y] \in I_1,
        \quad
        \fr{g}^{(n)} \in I_n,
    \end{equation*}
    which allows us to construct a chain,
    \begin{equation}
    \label{eq:cascade_derived_series}
        \fr{g} = I_0 \supseteq I_1 \supseteq \cdots \supseteq I_{k-1} \supseteq I_k = \{0\},
    \end{equation}
    where quotients of ideals, $I_i / I_{i+1}$, are all abelian by construction.
    \cref{eq:cascade_derived_series} is not necessarily the shortest sequence of abelian ideals,
    so abelian $k$-layer SSM has maximally $k$ derived length.
\end{proof}
\begin{proof}[Restricted SSM]
    It is trivial from \cref{eq:cascade_derived_series}.
    Each $I_i / I_{i+1}$ is constructed to be restricted,
    which by themselves are solvable with derived length $2$.
    Again, the series is not necessarily the shortest derived series,
    and restricted $k$-layer SSM has derived length at most $2k$.
\end{proof}
\subsection{Proof of \cref{thm:Kstacks}}
\label{app:sub_proofKstacks}
\begin{proof}
    Consider lifted Lie equation of $\mbf{S}_\fr{g}$ with a solvable matrix Lie algebra $\fr{g} \subset \fr{gl}(V)$:
    \begin{equation}
    \label{app:eq_kstack_solvable_Lie_eq}
        \dot{\mbf{G}}(t) = \mbf{A}_\fr{g}(x(t))\mbf{G}(t),
        \quad
        \mbf{G}(0) = I_V,
    \end{equation}
    where $\mbf{A}_{\fr{g}}(x) \in \fr{g}$.
    Let the derived length of $\fr{g}$ be $(k+1)$ for the notational clarity.

    Recall \cref{app:rem_solvable_k_tower}, for the derived series and short exact sequence of $\fr{g}$:
    \begin{equation}
    \label{app:eq_proofs_solvable_algebra_seq}
        \fr{g} = \fr{g}^{(0)} \supseteq \cdots \supseteq \fr{g}^{(k)} \supseteq \fr{g}^{(k+1)} = \{0\},
        \quad
        0 \to \fr{g}^{(k)} \xrightarrow{i} \fr{g}  \xrightarrow{\theta} \fr{g} / \fr{g}^{(k)} \to 0.
    \end{equation}
    Then, our interest is to decompose the integral curve $\mbf{G}(t)$ into curves integrating $\fr{g}^{(k)}$ and $\fr{g} / \fr{g}^{(k)}$.
    Let $G$ be a matrix Lie group locally integrating $\fr{g}$ near neighborhood of identity. Again locally, let $A_k$ be the normal subgroup integrating the ideal $\fr{g}^{(k)}$, and $G_k$ for the quotient Lie algebra $\fr{g} / \fr{g}^{(k)}$.
    We do not discuss global property of $(A_k, G, G_k)$, and shortly call them as local matrix Lie groups. 
    Restricted to sufficiently small neighborhoods,
    \begin{equation*}
        1 \to A_k \xrightarrow{\boldsymbol i} G \xrightarrow{\Theta} G_k \to 1,
    \end{equation*}
    where differentiation at the identity recovers \cref{app:eq_proofs_solvable_algebra_seq}.
    We follow common notational convention and do not distinguish
    $A_k$ or $\fr{g}^{(k)}$ from $\boldsymbol i(A_k)$ or $i(\fr{g}^{(k)})$,
    respectively.
    
    For the quotient part, its Lie equation is straightforward.
    We often silence variables without confusion for the clarify.
    With $\theta = d\Theta_e$:
    \begin{align}
        \mbf{G}_k &:= \Theta(\mbf{G}) \in G_k,
        \quad
        \bar{\mbf{A}} := \theta(\mbf{A}_\fr{g}) \in \fr{g} / \fr{g}^{(k)}, \nonumber
        \\
        \dot{\mbf{G}}_k
        &=
        \frac{d}{dt}\Theta(\mbf{G})
        =
        (d\Theta_{\mbf{G}})(\dot{\mbf{G}})
        =
        \theta(\mbf{A}_\fr{g})\Theta(\mbf{G})
        =
        \bar{\mbf{A}}\mbf{G}_k. \label{app:eq_quotient_LieEquation}
    \end{align}
    Now let us construct the Lie equation generated by $\fr{g}^{(k)}$. We take a smooth local section,
    \begin{equation*}
        \mbf{s}: G_k \to G,
        \quad
        \Theta(\mbf{s}(\mbf{G}_k(t))) = \mbf{G}_k(t).
    \end{equation*}
    As $\ker \Theta = \text{im} \: \boldsymbol i$ from the definition of the short exact sequence,
    we can characterize the abelian ideal.
    Given an inversion map on group $\iota: G \to G$,
    \begin{equation}
    \label{app:eq_abelian_ideal_LieEquation}
        \mbf{a}(t) := \mbf{Q}^{-1}(t)\mbf{G}(t),
        \quad
        \mbf{Q}(t) := \mbf{s}(\mbf{G}_k(t)) \in G,
        \quad
        \mbf{Q}^{-1}(t) = \iota(\mbf{Q}(t)).
    \end{equation}
    We can confirm $\mbf{a}(t)$ is indeed an element of the abelian ideal, sitting in $\ker \Theta$:
    \begin{equation*}
        \Theta(\mbf{Q}^{-1}\mbf{G}) = \Theta(\mbf{Q}^{-1})\Theta(\mbf{G}) = \Theta(\iota(\mbf{s}(\mbf{G}_k)))\mbf{G}_k = \mbf{G}_k^{-1}\mbf{G}_k = e_{G_k},
    \end{equation*}
    therefore any $\mbf{a}(t) \in A_k$.
    Differentiation gives us the abelian Lie equation of $\fr{g}^{(k)}$.
    \begin{equation}
    \label{app:eq_abelian_ideal_first_differentiation}
        \dot{\mbf{a}} = \left(\frac{d}{dt} \mbf{Q}^{-1} \right)\mbf{G} + \mbf{Q}^{-1}\dot{\mbf{G}}
        = -\mbf{Q}^{-1}\dot{\mbf{Q}}\mbf{Q}^{-1}\mbf{G} + \mbf{Q}^{-1}\dot{\mbf{G}}
        = (\mbf{Q}^{-1}\mbf{A}_\fr{g}\mbf{Q}-\mbf{Q}^{-1}\dot{\mbf{Q}})\mbf{a},
    \end{equation}
    where $\dot{\mbf{Q}}(t)$ follows,
    \begin{equation*}
        \dot{\mbf{Q}} = (d\mbf{s})_{\mbf{G}_k}(\dot{\mbf{G}}_k) = (d\mbf{s})_{\mbf{G}_k}(\bar{\mbf{A}}\mbf{G}_k).
    \end{equation*}
    Let us confirm $(\mbf{Q}^{-1}\mbf{A}_\fr{g}\mbf{Q}-\mbf{Q}^{-1}\dot{\mbf{Q}}) \in \fr{g}^{(k)}$
    to ensure \cref{app:eq_abelian_ideal_first_differentiation} is indeed the Lie equation of $\fr{g}^{(k)}$.
    
    First, note that $\mbf{Q}^{-1}\mbf{A}_\fr{g}\mbf{Q}$ is an adjoint action (or conjugation) of Lie group on its Lie algebra
    \cite{robbin2022introduction}.
    Being matrix Lie group and algebra, it is a standard matrix calculus:
    \begin{equation*}
        \theta(\mbf{Q}^{-1}\mbf{A}_\fr{g}\mbf{Q})
        =
        \theta(\text{Ad}_{\mbf{Q}^{-1}}\mbf{A}_\fr{g})
        =
        \text{Ad}_{\Theta(\mbf{Q}^{-1})}\bigl(\theta(\mbf{A}_\fr{g})\bigr)
        =
        \Theta(\mbf{Q}^{-1})\theta(\mbf{A}_\fr{g})\Theta(\mbf{Q})
        =
        \mbf{G}_k^{-1}\bar{\mbf{A}}\mbf{G}_k.
    \end{equation*}
    The second term $\mbf{Q}^{-1}\dot{\mbf{Q}}$ is simply derived from \cref{app:eq_quotient_LieEquation}:    
    \begin{equation*}
        \theta(\mbf{Q}^{-1}\dot{\mbf{Q}}) = \mbf{G}_k^{-1}\dot{\mbf{G}}_k = \mbf{G}_k^{-1}\bar{\mbf{A}}\mbf{G}_k,
    \end{equation*}
    then terms cancel out,
    \begin{equation*}
        \theta(\mbf{Q}^{-1}\mbf{A}_\fr{g}\mbf{Q}-\mbf{Q}^{-1}\dot{\mbf{Q}}) = 0.
    \end{equation*}
    Recall from \cref{app:eq_proofs_solvable_algebra_seq} that $\ker \theta = \text{im} \; i$.
    \cref{app:eq_abelian_ideal_first_differentiation} is an abelian Lie equation generated by the abelian ideal.
    
    Together, \cref{app:eq_kstack_solvable_Lie_eq} locally admits 2-layer cascade decomposition:
    \begin{align*}
        \mbf{G}(t) = \mbf{s}(\mbf{G}_k(t))\mbf{a}(t),
        \quad
        \dot{\mbf{G}}_k(t) &= \bar{\mbf{A}}(x(t))\mbf{G}_k(t),
        \quad
        \dot{\mbf{a}}(t) = \mbf{A}_k(x(t), \mbf{G}_k(t))\mbf{a}(t), \\
        \mbf{A}_k(x(t), \mbf{G}_k(t)) &= \mbf{Q}^{-1}(t)\mbf{A}_\fr{g}(x(t))\mbf{Q}(t)-\mbf{Q}^{-1}(t)\dot{\mbf{Q}}(t).
    \end{align*}
    As the base layer, $\fr{g} / \fr{g}^{(k)}$, is again solvable with derived length $k$,
    we can again faithfully represent $\fr{g} / \fr{g}^{(k)}$ as a matrix Lie algebra and iteratively apply the same decomposition near identity.
    Therefore, the lifted Lie equation of $\mbf{S}_\fr{g}$ decomposes into $(k+1)$-layer cascade of abelian Lie equations.
    With \cref{eq:flow_to_state_submersion}, decomposed Lie equations
    can locally simulate $\mbf{S}_\fr{g}$.

    Lastly, let us represent $(k+1)$-layer cascade of abelian Lie equations with abelian $(k+1)$-layer SSM.
    A Lie equation can be vectorized with the Kronecker product $\otimes$
    to be represented in SSM structure.
    In \cref{app:eq_kstack_solvable_Lie_eq},
    consider $\fr{g} \subset \fr{gl}(n,\bb{R})$.
    Let $\mbf{z} \in \bb{R}^{n^2}$ be a flattened $\mbf{G} \in \bb{R}^{n \times n}$:
    \begin{equation}
    \label{app:eq_Lie_to_SSM_vectorization}
        \dot{\mbf{z}}(t)
        =
        \bigl(
        I_n \otimes \mbf{A}_\fr{g}(x(t))
        \bigr)
        \mbf{z}(t),
        \quad
        \mbf{z}(0) = \text{vec}(I_n).
    \end{equation}
    Therefore, a solvable $\mbf{S}_\fr{g}$ can be decomposed and simulated by an abelian deep SSM.
\end{proof}
\cref{app:eq_Lie_to_SSM_vectorization} shows the price to pay in our construction.
    Lifting requires at most quadratic state dimension expansion.
    In general, $\dim \fr{g} \le n^2$.
    Especially for any finite-dimensional solvable Lie algebra $\fr{g}$ over $\bb{C}$,
    Lie's theorem states that $\fr{g}$ is faithfully represented by upper triangular matrices,
    therefore $\dim \fr{g} < n^2$.

\begin{remark}[Non-abelian extension]
\label{rem:non_abelian_extension_generalization}
    \cref{app:sub_proofKstacks} primarily aims the decomposition of solvable Lie algebras.
    The theorem constructs the basis for \cref{col:nilpotentization} and \cref{col:logdepth}.
    However, it is worth noting that the proof technique itself does not rely on the ideal $\fr{g}^{(k)}$ being abelian.
    The same local decomposition applies to general finite-dimensional Lie algebra extensions.
    
    While this depth-extension correspondence is theoretically appealing, we do not emphasize it in the main text.
    Many of the practical implications following \cref{thm:Kstacks}, including \cref{col:nilpotentization} and \cref{col:logdepth}, are already well-captured by the solvable Lie algebra decomposition.
    
    Note that \cref{thm:Kstacks} complements results in \citet{krener1977decomposition}.
    \citet{krener1977decomposition} proved \emph{split} extension admits cascade decomposition.
\end{remark}
\subsection{Proof of \cref{col:nilpotentization}}
\label{app:sub_proof_nilpotentization}
The flow of non-solvable system $\mbf{S}_\fr{g}$ demands infinite nontrivial terms in Magnus expansion.
Nilpotentization, or truncation, of the Magnus expansion is
an approximation technique that matches up to $c$-th order Magnus term
by class $c$ nilpotent Lie algebra \cite{iserles2008magnus,blanes2009magnus}.
The idea is to quotient $\fr{g}$ by its $c$-th order lower central series $\fr{g}^c$, as alluded in \cref{app:CDE}.
As $\fr{g}^c$ is an ideal in $\fr{g}$, the quotient $\fr{g}/\fr{g}^c$ constructs as a nilpotent Lie algebra in class $c$.

It is well known that
a nilpotent Lie algebra with class $c$ has maximal derived length of $\lceil \log_2 c \rceil + 1$ \cite{burde2012derived}.
Combined with \cref{thm:Kstacks}, abelian $k$-layer SSM can
express finite-dimensional nilpotent system up to class $(2^{k-1})$.
Then, the leading error term is $(2^{k-1}+1)$-th order Magnus term.
The magnitude of leading Magnus term $\|\Omega_{2^{k-1}+1}\| \lesssim \mcal{O}(\epsilon^{2^{k-1}+1})$ \cite{blanes2009magnus}.
Therefore, adding layers can exponentially increase the truncation order
and decrease the magnitude of leading error term.
\subsection{Proof of \cref{col:logdepth}}
\label{app:sub_proof_logdepth}
Here we formally introduce the word problem for a finite monoid and related concepts.
\begin{definition}[Word problem \cite{merrill2024illusion}]
\label{app:def_formal_wordproblem}
    Consider a finite alphabet set $\mcal{A}$.
    We denote the free monoid on $\mcal{A}$ as $\mcal{A}^\ast$, and call its elements as \emph{words}.
    Also consider a finite monoid $M$ generated by an image of alphabet map $\boldsymbol\varphi: \mcal{A} \to M$.
    Then, there exists a unique monoid homomorphism $\widehat{\boldsymbol\varphi}:\mcal{A}^\ast \to M$ such that:
    \begin{equation*}
        \widehat{\boldsymbol\varphi}(a) = \boldsymbol\varphi (a),
        \quad
        \widehat{\boldsymbol\varphi}(a_1a_2\ldots a_t) = \boldsymbol\varphi (a_1) \cdot \boldsymbol\varphi (a_2) \cdot \ldots \cdot \boldsymbol\varphi (a_t) \in M,
    \end{equation*}
    for alphabets $a, a_1, \cdots, a_t \in \mcal{A}$.
    The word problem for monoid is a decision problem such that for two arbitrary words $u, v \in \mcal{A}^\ast$ we decide whether $\widehat{\boldsymbol\varphi} (u) = \widehat{\boldsymbol\varphi} (v)$.
\end{definition}
As noted in \citet{merrill2024illusion}, the word problem in sequence modeling requires one to learn the word evaluation, $\widehat{\boldsymbol\varphi}$.
Learning $\widehat{\boldsymbol\varphi}$ is equivalent to expressing the correct state transition operation when a new input arrives.
The decision of $\widehat{\boldsymbol\varphi} (u) = \widehat{\boldsymbol\varphi} (v)$ then immediately follows. Without confusion, we often identify $\boldsymbol\varphi (a) \in M$ with an alphabet $a \in \mcal{A}$.
\begin{definition}[Lyndon words \cite{avdieiev2024affine}]
\label{app:def_LyndonWords}
    Let $\mcal{A}$ be totally ordered.
    For any words, we introduce the \emph{lexicographical} order on $\mcal{A}^*$:
    \begin{equation}
    a_1a_2\ldots a_i < b_1b_2\ldots b_j \quad \text{if} \begin{cases}
            a_1 = b_1, a_2 = b_2, \ldots, a_n = b_n, a_{n+1} < b_{n+1}\text{ for some } n \ge 0 \\
            \text{or} \\
            a_1 = b_1, \ldots, a_i = b_i \text{ and } i < j.
        \end{cases}
    \end{equation}
    Then, a word $w = a_1a_2\ldots a_i$ is \emph{Lyndon} if it is smaller than all of its cyclic permutations:
    \begin{equation*}
        a_1\ldots a_{n-1}a_n \ldots a_i < a_n\ldots a_i a_1\ldots a_{n-1}
        \quad
        \forall n \in \{2, \ldots, i\}.
    \end{equation*}
    and we call $a_1\ldots a_{n-1}$ and $a_n \ldots a_i$ as \emph{prefix} and \emph{suffix} of the word, respectively.
\end{definition}
\begin{proof}
Following \cref{app:def_formal_wordproblem}, 
the Chen-Fox-Lyndon theorem states that any words have a unique \emph{canonical} factorization in Lyndon words \cite{chen1958free,melanccon1989lyndon,avdieiev2024affine}.
Let us denote a set of Lyndon words on $\mcal{A}$ as $L$, and its subset up to length $T$ as $L_T$.
For any word $w$,
\begin{equation*}
    w = l_1l_2\ldots l_k,
    \quad
    l_1 \ge l_2 \cdots \ge l_k,
    \quad
    l_i \in L.
\end{equation*}
Notice that the canonical factorization of a Lyndon word is trivially itself.
Instead, Lyndon word has a \emph{costandard} factorization. For any Lyndon word $l$,
\begin{equation*}
    l = l_1l_2,
    \quad
    l_1, l_2 \in L,
\end{equation*}
where $l_2$ is the longest proper Lyndon suffix of $l$.

Clearly, the canonical factorization provides the key insight in our proof.
Intuitively, when word length is bounded at $T$,
it suffices to "memorize" Lyndon words and their composition up to length $T$ to solve the problem.

It is known that $L$ provides a basis of \emph{free} Lie algebra on $\mcal{A}$ \cite{avdieiev2024affine, walker2024log}.
Let $\fr{L}$ be a free Lie algebra generated by a finite set $\{e_a\}_{a \in \mcal{A}}$,
where each element is labeled by an alphabet.
The bracket rule is given:
\begin{equation}
    \text{b}[a] = e_a \in \fr{L} \text{ for } a \in \mcal{A},
    \quad
    \text{b}[l] = [\text{b}[l_1],\text{b}[l_2]],
\end{equation}
with $l=l_1l_2$ be the costandard factorization.
For a simple example, consider a Lyndon word $a_1a_2a_3$ with costandard factorization $(a_1, a_2a_3)$. Then, 
\begin{equation*}
    \text{b}[a_1a_2a_3] = [\text{b}[a_1], \text{b}[a_2a_3]] = [e_{a_1}, [\text{b}[a_2], \text{b}[a_3]]] = [e_{a_1}, [e_{a_2}, e_{a_3}]] \in \fr{L}.
\end{equation*}
Observe that Lyndon word length $3$ corresponds to the second-order Lie bracket.
Truncate the Lyndon words at $L_T$ implies that we consider every bracket of length $\ge T$ vanishes.
Effectively, we take a quotient of $\fr{L}$ with $\fr{L}^T$, its $T$-th order lower central series term.
As $\fr{L}^T$ is an ideal in $\fr{L}$, the quotient is also a Lie algebra.
Following, $\fr{g} = \fr{L} / \fr{L}^T$ is a finite-dimensional nilpotent Lie algebra with class $T$.

Now we invoke \cref{thm:Kstacks} just as we did in \cref{app:sub_proof_nilpotentization}.
For a class-$T$ nilpotent $\mbf{S}_\fr{g}$, we can construct abelian $\bigl(\lceil \log_2T \rceil + 1\bigr)$-layer SSM plus output layer to simulate flow \cite{burde2012derived}.
Each alphabets is then a piecewise-constant external input signal to $\mbf{S}_\fr{g}$.
\end{proof}
\subsection{Proof of \cref{col:logdepth_state}}
\label{app:sub_proof_logdepth_state}
The construction is natural from \cref{col:logdepth}.
Following logics and notations in \cref{app:sub_proof_logdepth},
we highlight that the number of Lyndon words determines the dimension of $\fr{g}$.
Let the number of alphabets $|\mcal{A}| = n$,
and denote $\fr{g}(T,n) = \fr{L} / \fr{L}^T$.
Then, \emph{Witt's formula} provides following \cite{magnus1966combinatorial, bouallegue2010primitive}:
\begin{equation*}
    \dim \fr{g}(T,n) = \sum_{m=1}^T \frac{1}{m} \sum_{d \mid m} n^d \ \mu\bigl(\frac{m}{d} \bigr),
\end{equation*}
where $\mu$ is the M\"obius function.